\section{Product of preschemes}
\label{section:I.3}

\subsection{Sum of preschemes.}
\label{subsection:I.3.1}

\phantomsection\egalabel{I.3.1.1}{3.1.1}
Let $(X_\alpha)$ be an arbitrary family of preschemes; let $X$ be a topological
space which is the \emph{sum} of the underlying spaces $X_\alpha$; $X$ is then
the union of pairwise disjoint open subspaces $X'_\alpha$, and for
each $\alpha$ there is a homeomorphism $\varphi_\alpha$ of $X_\alpha$ onto
$X'_\alpha$. If each of the $X'_\alpha$ is endowed with the sheaf
$(\varphi_\alpha)_*(\mathscr{O}_{X_\alpha})$, it is clear that $X$ becomes a
prescheme, which is called the \emph{sum} of the family of preschemes $(X_\alpha)$
and which is denoted by $\bigsqcup_\alpha X_\alpha$. If $Y$ is a prescheme,
the map $f\mapsto(f\circ\varphi_\alpha)$ is a \emph{bijection} of
the set $\operatorname{Hom}(X,Y)$ onto the product set
$\prod_\alpha\operatorname{Hom}(X_\alpha,Y)$. In particular, if the
$X_\alpha$ are $S$-preschemes, with structural morphisms $\psi_\alpha$,
$X$ is an $S$-prescheme by means of the unique morphism $\psi:X\to S$ such that
$\psi\circ\varphi_\alpha=\psi_\alpha$ for all $\alpha$. The sum of two
preschemes $X$, $Y$ is denoted by $X\sqcup Y$. It is immediate that if
$X=\operatorname{Spec}(A)$, $Y=\operatorname{Spec}(B)$, then $X\sqcup Y$
is canonically identified with $\operatorname{Spec}(A\times B)$.

\subsection{Product of preschemes.}
\label{subsection:I.3.2}

\begin{definition}[3.2.1]
\label{I.3.2.1}
Given two $S$-preschemes $X$, $Y$, we say that a triple
$(Z,p_1,p_2)$ consisting of an $S$-prescheme $Z$ and of two $S$-morphisms
$p_1:Z\to X$, $p_2:Z\to Y$ is a product of the
\oldpage[I]{105}
$S$-preschemes $X$ and $Y$ if, for every $S$-prescheme $T$, the map
$f\mapsto(p_1\circ f,p_2\circ f)$ is a bijection of the set of
$S$-morphisms from $T$ to $Z$ onto the set of pairs consisting of an
$S$-morphism $T\to X$ and an $S$-morphism $T\to Y$ (in other words, a
bijection
\[
  \operatorname{Hom}_S(T,Z)\xrightarrow{\sim}
  \operatorname{Hom}_S(T,X)\times\operatorname{Hom}_S(T,Y)).
\]
\end{definition}

This is therefore the general notion of the \emph{product} of two objects of a
category, applied to the category of $S$-preschemes (T, I, 1.1); in
particular, a product of two $S$-preschemes is \emph{unique} up to a unique
$S$-isomorphism. Because of this uniqueness, a product of two $S$-preschemes
$X$, $Y$ is most often denoted by
$X\times_S Y$ (or simply $X\times Y$ if no confusion is to be feared),
the morphisms $p_1$, $p_2$ (called the \emph{canonical projections} of
$X\times_S Y$ to $X$ and $Y$ respectively) being omitted from the notation.
If $g:T\to X$, $h:T\to Y$ are two $S$-morphisms, we shall denote by
$(g,h)_S$, or simply $(g,h)$, the $S$-morphism
$f:T\to X\times_S Y$ such that $p_1\circ f=g$, $p_2\circ f=h$. If $X'$, $Y'$
are two $S$-preschemes, $p'_1$, $p'_2$ the canonical projections of
$X'\times_S Y'$ (assumed to exist), and $u:X'\to X$, $v:Y'\to Y$ two
$S$-morphisms, we shall write $u\times_S v$ (or simply $u\times v$) for the
$S$-morphism $(u\circ p'_1,v\circ p'_2)_S$ from $X'\times_S Y'$ to
$X\times_S Y$.

When $S$ is an affine scheme with ring $A$, we often replace $S$ by $A$
in the preceding notation.

\begin{proposition}[3.2.2]
\label{I.3.2.2}
Let $X$, $Y$, $S$ be three affine schemes, and $B$, $C$, $A$ their respective
rings. Let $Z=\operatorname{Spec}(B\otimes_A C)$, and let $p_1$, $p_2$ be the
$S$-morphisms corresponding (\hyperref[I.2.2.4]{2.2.4}) to the
canonical $A$-homomorphisms $u:b\mapsto b\otimes1$ and
$v:c\mapsto1\otimes c$ of $B$ and $C$ into $B\otimes_A C$; then
$(Z,p_1,p_2)$ is a product of $X$ and $Y$.
\end{proposition}

By virtue of (\hyperref[I.2.2.4]{2.2.4}), everything comes down to verifying
that if, to every $A$-homomorphism $f:B\otimes_A C\to L$ (where $L$ is an
$A$-algebra), we associate the pair $(f\circ u,f\circ v)$, we define a bijection
\[
  \operatorname{Hom}_A(B\otimes_A C,L)\xrightarrow{\sim}
  \operatorname{Hom}_A(B,L)\times\operatorname{Hom}_A(C,L)\footnotemark,
\]
\footnotetext{The notation $\operatorname{Hom}_A$ here denotes the set of
homomorphisms of $A$-algebras.} which follows immediately from the definitions
and from the relation $b\otimes c=(b\otimes1)(1\otimes c)$.

\begin{corollary}[3.2.3]
\label{I.3.2.3}
Let $T$ be an affine scheme with ring $D$, and
$\alpha=({}^a\rho,\widetilde{\rho})$
(resp. $\beta=({}^a\sigma,\widetilde{\sigma})$) an $S$-morphism $T\to X$
(resp. $T\to Y$), where $\rho$ (resp. $\sigma$) is an $A$-homomorphism of $B$
(resp. $C$) into $D$; then
$(\alpha,\beta)_S=({}^a\tau,\widetilde{\tau})$, where $\tau$ is the homomorphism
$B\otimes_A C\to D$ such that $\tau(b\otimes c)=\rho(b)\sigma(c)$.
\end{corollary}

\begin{proposition}[3.2.4]
\label{I.3.2.4}
Let $f:S'\to S$ be a monomorphism of preschemes (T, I, 1.1), and $X$, $Y$ two
$S'$-preschemes, which are also regarded as $S$-preschemes by means of
$f$. Then every product of the $S$-preschemes $X$, $Y$ is a product of the
$S'$-preschemes $X$, $Y$, and conversely.
\end{proposition}

Let $\varphi:X\to S'$, $\psi:Y\to S'$ be the structural morphisms. If $T$
is an $S$-prescheme and $u:T\to X$, $v:T\to Y$ are two $S$-morphisms, we have by
definition $f\circ\varphi\circ u=f\circ\psi\circ v=\theta$, the structural
morphism of $T$; the hypothesis on $f$ implies
$\varphi\circ u=\psi\circ v=\theta'$, and we see that we can regard $T$
as an $S'$-prescheme with structural morphism $\theta'$, and $u$ and $v$ as
$S'$-morphisms. The conclusion of the proposition follows immediately,
taking into account (\hyperref[I.3.2.1]{3.2.1}).

\begin{corollary}[3.2.5]
\label{I.3.2.5}
Let $X$, $Y$ be two $S$-preschemes, $\varphi:X\to S$, $\psi:Y\to S$ their
structural morphisms, and $S'$ an open subset of $S$ such that
$\varphi(X)\subset S'$, $\psi(Y)\subset S'$. Every product of the $S$-preschemes
$X$, $Y$ is also a product of the $S'$-preschemes $X$, $Y$, and conversely.
\end{corollary}

\oldpage[I]{106}
It suffices to apply (\hyperref[I.3.2.4]{3.2.4}) to the canonical injection
$S'\to S$.

\begin{theorem}[3.2.6]
\label{I.3.2.6}
Given two $S$-preschemes $X$, $Y$, there exists a product
$X\times_S Y$.
\end{theorem}

We shall proceed in several steps.

\begin{lemma}[3.2.6.1]
\label{I.3.2.6.1}
Let $(Z,p,q)$ be a product of $X$ and $Y$, and $U$, $V$ open subsets of
$X$, $Y$ respectively. If we set
$W=p^{-1}(U)\cap q^{-1}(V)$, the triple consisting of $W$ and of the restrictions
of $p$ and $q$ to $W$ \textup{(regarded as morphisms $W\to U$,
$W\to V$ respectively)} is a product of $U$ and $V$.
\end{lemma}

Indeed, if $T$ is an $S$-prescheme, we can identify the $S$-morphisms
$T\to W$ with the $S$-morphisms $T\to Z$ mapping $T$ into $W$. If then
$g:T\to U$, $h:T\to V$ are two arbitrary $S$-morphisms, we can
regard them as $S$-morphisms of $T$ into $X$ and $Y$ respectively, and
by hypothesis there is therefore one and only one $S$-morphism $f:T\to Z$ such
that $g=p\circ f$, $h=q\circ f$. Since $p(f(T))\subset U$, $q(f(T))\subset V$,
we have
\[
  f(T)\subset p^{-1}(U)\cap q^{-1}(V)=W,
\]
whence our assertion.

\begin{lemma}[3.2.6.2]
\label{I.3.2.6.2}
Let $Z$ be an $S$-prescheme, $p:Z\to X$, $q:Z\to Y$ two $S$-morphisms,
$(U_\alpha)$ an open covering of $X$, and $(V_\lambda)$ an open covering
of $Y$. Suppose that for every pair $(\alpha,\lambda)$, the
$S$-prescheme
$W_{\alpha\lambda}=p^{-1}(U_\alpha)\cap q^{-1}(V_\lambda)$ and the
restrictions of $p$ and $q$ to $W_{\alpha\lambda}$ constitute a product of
$U_\alpha$ and $V_\lambda$. Then $(Z,p,q)$ is a product of $X$ and $Y$.
\end{lemma}

Let us first show that, if $f_1$, $f_2$ are two $S$-morphisms $T\to Z$, the
relations $p\circ f_1=p\circ f_2$ and $q\circ f_1=q\circ f_2$ imply
$f_1=f_2$. Indeed, $Z$ is the union of the $W_{\alpha\lambda}$, so the
$f_1^{-1}(W_{\alpha\lambda})$ form an open covering of $T$, and
the same holds for the $f_2^{-1}(W_{\alpha\lambda})$. Moreover, we have
\[
\begin{aligned}
  f_1^{-1}(W_{\alpha\lambda})
  &=f_1^{-1}(p^{-1}(U_\alpha))\cap
    f_1^{-1}(q^{-1}(V_\lambda)) \\
  &=f_2^{-1}(p^{-1}(U_\alpha))\cap
    f_2^{-1}(q^{-1}(V_\lambda))
   =f_2^{-1}(W_{\alpha\lambda})
\end{aligned}
\]
by hypothesis, and everything comes down to seeing that the restrictions of
$f_1$ and $f_2$ to $f_1^{-1}(W_{\alpha\lambda})=f_2^{-1}(W_{\alpha\lambda})$ are
identical for every pair of indices. But since these restrictions can be
regarded as $S$-morphisms of
$f_1^{-1}(W_{\alpha\lambda})$ into $W_{\alpha\lambda}$, our assertion
follows from the hypotheses and from Definition (\hyperref[I.3.2.1]{3.2.1}).

Suppose now that two $S$-morphisms $g:T\to X$, $h:T\to Y$ are given.
Set
$T_{\alpha\lambda}=g^{-1}(U_\alpha)\cap h^{-1}(V_\lambda)$; the
$T_{\alpha\lambda}$ form an open covering of $T$. By hypothesis,
there exists an $S$-morphism $f_{\alpha\lambda}$ such that
$p\circ f_{\alpha\lambda}$ and $q\circ f_{\alpha\lambda}$ are the
respective restrictions of $g$ and $h$ to $T_{\alpha\lambda}$. Moreover,
let us show that the restrictions of $f_{\alpha\lambda}$ and
$f_{\beta\mu}$ to
$T_{\alpha\lambda}\cap T_{\beta\mu}$ coincide, which will complete the
proof of (\hyperref[I.3.2.6.2]{3.2.6.2}). Now, the images of
$T_{\alpha\lambda}\cap T_{\beta\mu}$ under $f_{\alpha\lambda}$ and under
$f_{\beta\mu}$ are contained in
$W_{\alpha\lambda}\cap W_{\beta\mu}$ by definition. Since
\[
  W_{\alpha\lambda}\cap W_{\beta\mu}
  =p^{-1}(U_\alpha\cap U_\beta)\cap
   q^{-1}(V_\lambda\cap V_\mu),
\]
it follows from (\hyperref[I.3.2.6.1]{3.2.6.1}) that
$W_{\alpha\lambda}\cap W_{\beta\mu}$ and the restrictions of $p$ and $q$ to
this prescheme constitute a product of
$U_\alpha\cap U_\beta$ and $V_\lambda\cap V_\mu$. Since
$p\circ f_{\alpha\lambda}$ and $p\circ f_{\beta\mu}$ coincide in
$T_{\alpha\lambda}\cap T_{\beta\mu}$ and the same holds for
$q\circ f_{\alpha\lambda}$ and $q\circ f_{\beta\mu}$, we see that
$f_{\alpha\lambda}$ and $f_{\beta\mu}$ coincide in
$T_{\alpha\lambda}\cap T_{\beta\mu}$, Q.E.D.

\oldpage[I]{107}
\begin{lemma}[3.2.6.3]
\label{I.3.2.6.3}
Let $(U_\alpha)$ be an open covering of $X$, $(V_\lambda)$ an open
covering of $Y$, and suppose that for every pair
$(\alpha,\lambda)$ there exists a product of $U_\alpha$ and $V_\lambda$;
then there exists a product of $X$ and $Y$.
\end{lemma}

Applying Lemma (\hyperref[I.3.2.6.1]{3.2.6.1}) to the open sets
$U_\alpha\cap U_\beta$ and $V_\lambda\cap V_\mu$, we see that there exists a
product of the $S$-preschemes induced by $X$ and $Y$ respectively on these
open sets; moreover, the uniqueness of the product shows that, if we set
$i=(\alpha,\lambda)$, $j=(\beta,\mu)$, there is a canonical isomorphism
$h_{ij}$ (resp. $h_{ji}$) of this product onto an $S$-prescheme $W_{ij}$
(resp. $W_{ji}$) induced by $U_\alpha\times_S V_\lambda$
(resp. $U_\beta\times_S V_\mu$) on an open set;
$f_{ij}=h_{ij}\circ h_{ji}^{-1}$ is therefore an isomorphism of $W_{ji}$ onto
$W_{ij}$. Moreover, for a third pair $k=(\gamma,\nu)$, we have
$f_{ik}=f_{ij}\circ f_{jk}$ in $W_{ki}\cap W_{kj}$, as follows from
applying (\hyperref[I.3.2.6.1]{3.2.6.1}) to the open sets
$U_\alpha\cap U_\beta\cap U_\gamma$ and
$V_\lambda\cap V_\mu\cap V_\nu$, in $U_\beta$ and $V_\mu$
respectively. Consequently there exist a prescheme $Z$, an open covering
$(Z_i)$ of the underlying space of $Z$ and, for each $i$, an isomorphism $g_i$
of the induced prescheme $Z_i$ onto the prescheme
$U_\alpha\times_S V_\lambda$, such that for every pair $(i,j)$ we have
$f_{ij}=g_i\circ g_j^{-1}$
(\hyperref[I.2.3.1]{2.3.1}); moreover, we have
$g_i(Z_i\cap Z_j)=W_{ij}$. If $p_i$, $q_i$, $\theta_i$ are the projections
and the structural morphism of the $S$-prescheme
$U_\alpha\times_S V_\lambda$, one verifies at once that
$p_i\circ g_i=p_j\circ g_j$ in $Z_i\cap Z_j$, and likewise for the two
other morphisms. We can therefore define morphisms of preschemes
$p:Z\to X$ (resp. $q:Z\to Y$, $\theta:Z\to S$) by the condition that $p$
(resp. $q$, $\theta$) coincides with $p_i\circ g_i$
(resp. $q_i\circ g_i$, $\theta_i\circ g_i$) in each of the $Z_i$; $Z$,
endowed with $\theta$, is then an $S$-prescheme. Let us now show that
$Z'_i=p^{-1}(U_\alpha)\cap q^{-1}(V_\lambda)$ is equal to $Z_i$. For every
index $j=(\beta,\mu)$, we have
$Z_j\cap Z'_i=g_j^{-1}(p_j^{-1}(U_\alpha)\cap q_j^{-1}(V_\lambda))$. Now,
\[
  p_j^{-1}(U_\alpha)\cap q_j^{-1}(V_\lambda)
  =p_j^{-1}(U_\alpha\cap U_\beta)\cap
   q_j^{-1}(V_\lambda\cap V_\mu)~;
\]
by virtue of (\hyperref[I.3.2.6.1]{3.2.6.1}), the restrictions of $p_j$ and
$q_j$ to
$p_j^{-1}(U_\alpha)\cap q_j^{-1}(V_\lambda)$ define on this
$S$-prescheme a structure of product of $U_\alpha\cap U_\beta$ and
$V_\lambda\cap V_\mu$; but the uniqueness of the product then implies that
$p_j^{-1}(U_\alpha)\cap q_j^{-1}(V_\lambda)=W_{ji}$. Consequently we have
$Z_j\cap Z'_i=Z_j\cap Z_i$ for all $j$, whence $Z'_i=Z_i$. We then deduce
from (\hyperref[I.3.2.6.2]{3.2.6.2}) that $(Z,p,q)$ is a product of $X$ and
$Y$.

\begin{lemma}[3.2.6.4]
\label{I.3.2.6.4}
Let $\varphi:X\to S$, $\psi:Y\to S$ be the structural morphisms of $X$
and $Y$, $(S_i)$ an open covering of $S$, and set
$X_i=\varphi^{-1}(S_i)$, $Y_i=\psi^{-1}(S_i)$. If each of the products
$X_i\times_S Y_i$ exists, then $X\times_S Y$ exists.
\end{lemma}

By (\hyperref[I.3.2.6.3]{3.2.6.3}), everything comes down to proving that the
products $X_i\times_S Y_j$ exist for all $i$ and $j$. Set
$X_{ij}=X_i\cap X_j=\varphi^{-1}(S_i\cap S_j)$,
$Y_{ij}=Y_i\cap Y_j=\psi^{-1}(S_i\cap S_j)$; by virtue of
(\hyperref[I.3.2.6.1]{3.2.6.1}), the product
$Z_{ij}=X_{ij}\times_S Y_{ij}$ exists. Now note that if $T$ is an
$S$-prescheme and $g:T\to X_i$, $h:T\to Y_j$ are $S$-morphisms, we
necessarily have
$\varphi(g(T))=\psi(h(T))\subset S_i\cap S_j$, by the definition of an
$S$-morphism, hence $g(T)\subset X_{ij}$ and $h(T)\subset Y_{ij}$; it is
then immediate that $Z_{ij}$ is a product of $X_i$ and $Y_j$.

\begin{env}[3.2.6.5]
\label{I.3.2.6.5}
We can now complete the proof of Theorem
(\hyperref[I.3.2.6]{3.2.6}). If $S$ is an \emph{affine scheme}, there are
coverings $(U_\alpha)$, $(V_\lambda)$ of $X$ and $Y$ respectively,
consisting of affine open sets; since $U_\alpha\times_S V_\lambda$ exists by
virtue of (\hyperref[I.3.2.2]{3.2.2}), the same holds for
$X\times_S Y$ by (\hyperref[I.3.2.6.3]{3.2.6.3}). If $S$ is an
arbitrary prescheme, there is a covering $(S_i)$ of $S$ consisting of affine
open sets. If $\varphi:X\to S$, $\psi:Y\to S$ are the structural
morphisms, and if we set $X_i=\varphi^{-1}(S_i)$,
$Y_i=\psi^{-1}(S_i)$, the products $X_i\times_{S_i}Y_i$ exist by the
\oldpage[I]{108}
foregoing; but then the products $X_i\times_S Y_i$ also exist
(\hyperref[I.3.2.5]{3.2.5}), hence the same holds for $X\times_S Y$ by
(\hyperref[I.3.2.6.4]{3.2.6.4}).
\end{env}

\begin{corollary}[3.2.7]
\label{I.3.2.7}
Let $Z=X\times_S Y$ be the product of two $S$-preschemes, $p$, $q$ the
projections of $Z$ to $X$ and $Y$, and $\varphi$ (resp. $\psi$) the structural
morphism of $X$ (resp. $Y$). Let $S'$ be an open subset of $S$, and $U$
(resp. $V$) an open subset of $X$ (resp. $Y$) contained in
$\varphi^{-1}(S')$ (resp. $\psi^{-1}(S')$). Then the product
$U\times_{S'}V$ is canonically identified with the prescheme induced by $Z$ on
$p^{-1}(U)\cap q^{-1}(V)$ (regarded as an $S'$-prescheme). Moreover, if
$f:T\to X$, $g:T\to Y$ are $S$-morphisms such that $f(T)\subset U$,
$g(T)\subset V$, the $S'$-morphism $(f,g)_{S'}$ is identified with the
restriction of $(f,g)_S$ to $p^{-1}(U)\cap q^{-1}(V)$.
\end{corollary}

This follows from (\hyperref[I.3.2.5]{3.2.5}) and
(\hyperref[I.3.2.6.1]{3.2.6.1}).

\begin{env}[3.2.8]
\label{I.3.2.8}
Let $(X_\alpha)$, $(Y_\lambda)$ be two families of $S$-preschemes, and $X$
(resp. $Y$) the sum of the family $(X_\alpha)$ (resp. $(Y_\lambda)$)
(\hyperref[subsection:I.3.1]{3.1}). Then $X\times_S Y$ is identified with the
\emph{sum} of the family $(X_\alpha\times_S Y_\lambda)$; this follows
at once from (\hyperref[I.3.2.6.3]{3.2.6.3}).
\end{env}

\subsection{Formal properties of the product; change of base prescheme.}
\label{subsection:I.3.3}

\begin{env}[3.3.1]
\label{I.3.3.1}
The reader will note that all the properties stated in this section,
except (\hyperref[I.3.3.13]{3.3.13}), (\hyperref[I.3.3.15]{3.3.15}) and the assertion about sums in the last sentence of
(\hyperref[I.3.3.10]{3.3.10}), are valid without modification in
\emph{every category}, whenever the products which occur in
the statements \emph{exist} (for it is clear that the notions of $S$-object and
of $S$-morphism can be defined exactly as in
(\hyperref[subsection:I.2.5]{2.5}) for every object $S$ of the category).
\end{env}

\begin{env}[3.3.2]
\label{I.3.3.2}
In the first place, $X\times_S Y$ is a \emph{covariant bifunctor} in $X$ and
$Y$ on the category of $S$-preschemes: indeed, it suffices to note
that the diagram
\[
  \xymatrix{
    X\times Y\ar[r]^{f\times 1}\ar[d] &
    X'\times Y\ar[r]^{f'\times 1}\ar[d] &
    X''\times Y\ar[d]\\
    X\ar[r]_f & X'\ar[r]_{f'} & X''
  }
\]
is commutative.
\end{env}

\begin{proposition}[3.3.3]
\label{I.3.3.3}
For every $S$-prescheme $X$, the first (resp. second) projection of
$X\times_S S$ (resp. $S\times_S X$) is a functorial isomorphism of
$X\times_S S$ (resp. $S\times_S X$) onto $X$, whose inverse isomorphism
is $(1_X,\varphi)_S$ (resp. $(\varphi,1_X)_S$), where $\varphi$ denotes
the structural morphism $X\to S$; we can therefore write, up to canonical
isomorphism,
\[
  X\times_S S=S\times_S X=X.
\]
\end{proposition}

It suffices to prove that the triple $(X,1_X,\varphi)$ is a product of $X$
and $S$. Now, if $T$ is an $S$-prescheme, the only $S$-morphism of $T$ into
$S$ is necessarily the structural morphism $\psi:T\to S$. If $f$ is
an $S$-morphism of $T$ into $X$, we necessarily have
$\psi=\varphi\circ f$, whence our assertion.

\begin{corollary}[3.3.4]
\label{I.3.3.4}
Let $X$, $Y$ be two $S$-preschemes, and $\varphi:X\to S$, $\psi:Y\to S$ their
structural morphisms. If $X$ is canonically identified with $X\times_S S$
and $Y$ with $S\times_S Y$, the projections $X\times_S Y\to X$ and
$X\times_S Y\to Y$ are identified with $1_X\times\psi$ and
$\varphi\times 1_Y$ respectively.
\end{corollary}

The verification is immediate and is left to the reader.

\begin{env}[3.3.5]
\label{I.3.3.5}
One can define, in the same way as in
(\hyperref[subsection:I.3.2]{3.2}), the product of an
\oldpage[I]{109}
arbitrary finite number $n$ of $S$-preschemes; the existence of these products
follows from (\hyperref[I.3.2.6]{3.2.6}) by induction on $n$, noting
that
$(X_1\times_S X_2\times_S\cdots\times_S X_{n-1})\times_S X_n$ satisfies
the definition of the product. The uniqueness of the product implies, as in
every category, its \emph{commutativity} and \emph{associativity} properties.
If, for example, $p_1$, $p_2$, $p_3$ denote the projections of
$X_1\times_S X_2\times_S X_3$, and if this prescheme is identified with
$(X_1\times_S X_2)\times_S X_3$, the projection to $X_1\times_S X_2$ is
identified with $(p_1,p_2)_S$.
\end{env}

\begin{env}[3.3.6]
\label{I.3.3.6}
Let $S$, $S'$ be two preschemes, and $\varphi:S'\to S$ a morphism, which makes
$S'$ an $S$-prescheme. For every $S$-prescheme $X$, consider the product
$X\times_S S'$, and let $p$ and $\pi'$ be its projections to $X$ and $S'$
respectively. Endowed with $\pi'$, this product is an $S'$-prescheme; when it
is regarded as such, it is denoted by $X_{(S')}$ or $X_{(\varphi)}$, and it is
said to be the prescheme obtained by \emph{extension of the base prescheme} from
$S$ to $S'$, by means of the morphism $\varphi$, or the \emph{inverse image} of
$X$ under $\varphi$. Note that if $\pi$ is the structural morphism of $X$ and
$\theta$ the structural morphism of $X\times_S S'$ regarded as an
$S$-prescheme, the diagram
\phantomsection\label{I.3.3.6.base-change-diagram}
\[
  \xymatrix{
    X\ar[d]_{\pi} &
    X_{(S')}\ar[l]_{p}\ar[ld]_{\theta}\ar[d]^{\pi'}\\
    S & S'\ar[l]^{\varphi}
  }
\]
is commutative.
\end{env}

\begin{env}[3.3.7]
\label{I.3.3.7}
With the notation of (\hyperref[I.3.3.6]{3.3.6}), for every
$S$-morphism $f:X\to Y$, we also denote by $f_{(S')}$ the $S'$-morphism
$f\times_S 1:X_{(S')}\to Y_{(S')}$, and we say that $f_{(S')}$ is
the \emph{inverse image} of the morphism $f$ under $\varphi$. $X_{(S')}$ is
therefore a \emph{covariant functor} in $X$, from the category of
$S$-preschemes to that of $S'$-preschemes.
\end{env}

\begin{env}[3.3.8]
\label{I.3.3.8}
The prescheme $X_{(S')}$ can also be regarded as the solution of a
\emph{universal mapping problem}: every $S'$-prescheme $T$ is also
an $S$-prescheme by means of $\varphi$; every $S$-morphism $g:T\to X$
can then be written in one and only one way as $g=p\circ f$, where $f$ is an
$S'$-morphism $T\to X_{(S')}$, as follows from the definition of the product
applied to the $S$-morphisms $g$ and $\psi:T\to S'$ (the structural morphism
of $T$).
\end{env}

\begin{proposition}[3.3.9]
\label{I.3.3.9}
\textup{(``transitivity of the extension of base preschemes'').} Let
$S''$ be a prescheme and $\varphi':S''\to S'$ a morphism. For every
$S$-prescheme $X$, there exists a canonical functorial isomorphism of the
$S''$-prescheme $(X_{(\varphi)})_{(\varphi')}$ onto the $S''$-prescheme
$X_{(\varphi\circ\varphi')}$.
\end{proposition}

Indeed, let $T$ be an $S''$-prescheme, $\psi$ its structural morphism, and
$g$ an $S$-morphism of $T$ into $X$ ($T$ being regarded as an
$S$-prescheme with structural morphism $\varphi\circ\varphi'\circ\psi$).
Since $T$ is also an $S'$-prescheme with structural morphism
$\varphi'\circ\psi$, we can write $g=p\circ g'$, where $g'$ is an
$S'$-morphism $T\to X_{(\varphi)}$, and then $g'=p'\circ g''$, where $g''$ is
an $S''$-morphism $T\to(X_{(\varphi)})_{(\varphi')}$:
\phantomsection\label{I.3.3.9.transitivity-diagram}
\[
  \xymatrix{
    X\ar[d]_{\pi} &
    X_{(\varphi)}\ar[l]_{p}\ar[d]_{\pi'} &
    (X_{(\varphi)})_{(\varphi')}\ar[l]_{p'}\ar[d]_{\pi''}\\
    S & S'\ar[l]^{\varphi} & S''\ar[l]^{\varphi'}
  }
\]
\oldpage[I]{110}
Whence the proposition, because of the uniqueness of the solution of a
universal mapping problem.

This result can also be expressed by writing the equality (up to canonical
isomorphism) $(X_{(S')})_{(S'')}=X_{(S'')}$, if no confusion is to be
feared, or also
\begin{equation}
  (X\times_S S')\times_{S'}S''=X\times_S S''~;
  \tag{3.3.9.1}\label{I.3.3.9.1}
\end{equation}
the functorial character of the isomorphism defined in
(\hyperref[I.3.3.9]{3.3.9}) is likewise expressed by the formula of
\emph{transitivity of inverse images of morphisms}
\begin{equation}
  (f_{(S')})_{(S'')}=f_{(S'')}
  \tag{3.3.9.2}\label{I.3.3.9.2}
\end{equation}
for every $S$-morphism $f:X\to Y$.

\begin{corollary}[3.3.10]
\label{I.3.3.10}
If $X$ and $Y$ are two $S$-preschemes, there exists a canonical functorial
isomorphism of the $S'$-prescheme $X_{(S')}\times_{S'}Y_{(S')}$ onto the
$S'$-prescheme $(X\times_S Y)_{(S')}$.
\end{corollary}

Indeed, we have, up to canonical isomorphisms,
\[
  (X\times_S S')\times_{S'}(Y\times_S S')
  =X\times_S(Y\times_S S')=(X\times_S Y)\times_S S'
\]
by virtue of (\hyperref[I.3.3.9.1]{3.3.9.1}) and of the associativity of
products of $S$-preschemes.

The functorial character of the isomorphism defined in
(\hyperref[I.3.3.10]{3.3.10}) is expressed by the formula
\begin{equation}
  (u_{(S')},v_{(S')})_{S'}=((u,v)_S)_{(S')}
  \tag{3.3.10.1}\label{I.3.3.10.1}
\end{equation}
for every pair of $S$-morphisms $u:T\to X$, $v:T\to Y$.

In other words, the inverse image functor $X_{(S')}$
\emph{commutes with the formation of products}; it also commutes with the
formation of sums (\hyperref[I.3.2.8]{3.2.8}).

\begin{corollary}[3.3.11]
\label{I.3.3.11}
Let $Y$ be an $S$-prescheme, and $f:X\to Y$ a morphism which makes $X$ a
$Y$-prescheme (and consequently also an $S$-prescheme). The prescheme
$X_{(S')}$ is then identified with the product $X\times_Y Y_{(S')}$, the
projection $X\times_Y Y_{(S')}\to Y_{(S')}$ being identified with $f_{(S')}$.
\end{corollary}

Let $\psi:Y\to S$ be the structural morphism of $Y$; we have the commutative
diagram
\phantomsection\label{I.3.3.11.diagram}
\[
  \xymatrix{
    S'\ar[d] &
    Y_{(S')}\ar[l]_{\psi_{(S')}}\ar[d] &
    X_{(S')}\ar[l]_{f_{(S')}}\ar[d]\\
    S & Y\ar[l]^{\psi} & X\ar[l]^{f}
  }
\]
Now $Y_{(S')}$ is identified with $S'_{(\psi)}$, and $X_{(S')}$ with
$S'_{(\psi\circ f)}$; taking into account (\hyperref[I.3.3.9]{3.3.9}) and
(\hyperref[I.3.3.4]{3.3.4}), the corollary follows.

\begin{env}[3.3.12]
\label{I.3.3.12}
Let $f:X\to X'$, $g:Y\to Y'$ be two $S$-morphisms which are
\emph{monomorphisms} of preschemes (T, I, 1.1); then $f\times_S g$ is a
\emph{monomorphism}. Indeed, if $p$ and $q$ are the projections of
$X\times_S Y$, $p'$, $q'$ those of $X'\times_S Y'$, and $u$, $v$ two
$S$-morphisms $T\to X\times_S Y$, the relation
$(f\times_S g)\circ u=(f\times_S g)\circ v$ implies
$p'\circ(f\times_S g)\circ u=p'\circ(f\times_S g)\circ v$, in other words
$f\circ p\circ u=f\circ p\circ v$, and since $f$ is a monomorphism,
$p\circ u=p\circ v$; using the fact that $g$ is a monomorphism, we
likewise obtain $q\circ u=q\circ v$, whence $u=v$.
\oldpage[I]{111}
It follows that for every extension $S'\to S$ of the base prescheme,
\phantomsection\label{I.3.3.12.base-change-monomorphism}
\[
  f_{(S')}:X_{(S')}\to X'_{(S')}
\]
is a monomorphism.
\end{env}

\begin{env}[3.3.13]
\label{I.3.3.13}
Let $S$, $S'$ be two affine schemes with respective rings $A$, $A'$; a
morphism $S'\to S$ therefore corresponds to a ring homomorphism $A\to A'$.
If $X$ is an $S$-prescheme, we shall then also denote by $X_{(A')}$ or
$X\otimes_A A'$ the $S'$-prescheme $X_{(S')}$; when $X$ is itself
affine with ring $B$, $X_{(A')}$ is affine with ring
$B_{(A')}=B\otimes_A A'$, obtained by extending to $A'$ the ring of scalars
of the $A$-algebra $B$.
\end{env}

\begin{env}[3.3.14]
\label{I.3.3.14}
With the notation of (\hyperref[I.3.3.6]{3.3.6}), for every
\emph{$S$-morphism} $f:S'\to X$, $f'=(f,1_{S'})_S$ is an $S'$-morphism
$S'\to X'=X_{(S')}$ such that $p\circ f'=f$, $\pi'\circ f'=1_{S'}$, in other
words an \emph{$S'$-section of $X'$}; and conversely, if $f'$ is such an
$S'$-section, $f=p\circ f'$ is an $S$-morphism $S'\to X$. We thus define
a canonical \emph{one-to-one correspondence}
\phantomsection\label{I.3.3.14.correspondence}
\[
  \operatorname{Hom}_S(S',X)\overset{\sim}{\longrightarrow}
  \operatorname{Hom}_{S'}(S',X').
\]
We say that $f'$ is the \emph{graph morphism} of $f$, and we denote it by
$\Gamma_f$.
\end{env}

\begin{env}[3.3.15]
\label{I.3.3.15}
Given a prescheme $X$, which can always be regarded as a
$\mathbf{Z}$-prescheme, it follows in particular from
(\hyperref[I.3.3.14]{3.3.14}) that the $X$-sections of
$X\otimes_{\mathbf{Z}}\mathbf{Z}[T]$ (where $T$ is an indeterminate)
correspond one-to-one to the \emph{morphisms}
$X\to\mathbf{Z}[T]$. Let us show that these $X$-sections also correspond
one-to-one to the \emph{sections of the structure sheaf
$\mathcal{O}_X$ over $X$}. Indeed, let $(U_\alpha)$ be a covering
of $X$ by affine open sets; let
$u:X\to X\otimes_{\mathbf{Z}}\mathbf{Z}[T]$ be an $X$-morphism and let
$u_\alpha$ be its restriction to $U_\alpha$; if $A_\alpha$ is the ring of the
affine scheme $U_\alpha$, $U_\alpha\otimes_{\mathbf{Z}}\mathbf{Z}[T]$ is an
affine scheme with ring $A_\alpha[T]$ (\hyperref[I.3.2.2]{3.2.2}), and
$u_\alpha$ corresponds canonically to an $A_\alpha$-homomorphism
$A_\alpha[T]\to A_\alpha$ (1.7.3). Now, such a homomorphism is completely
determined by giving the image of $T$ in $A_\alpha$, say
$s_\alpha\in A_\alpha=\Gamma(U_\alpha,\mathcal{O}_X)$, and if we express that
the restrictions of $u_\alpha$ and of $u_\beta$ to an affine open set
$V\subset U_\alpha\cap U_\beta$ coincide, we see at once that $s_\alpha$
and $s_\beta$ coincide in $V$; hence the family $(s_\alpha)$ consists of the
restrictions to the $U_\alpha$ of a section $s$ of $\mathcal{O}_X$ over
$X$; and conversely, it is clear that such a section defines a family
$(u_\alpha)$ of morphisms which are the restrictions to the $U_\alpha$ of an
$X$-morphism $X\to X\otimes_{\mathbf{Z}}\mathbf{Z}[T]$. This result will be
generalized in II, 1.7.12.
\end{env}

\subsection{Points of a prescheme with values in a prescheme; geometric
points}
\label{subsection:I.3.4}

\begin{env}[3.4.1]
\label{I.3.4.1}
Let $X$ be a prescheme; for every prescheme $T$, we also denote by $X(T)$
the set $\operatorname{Hom}(T,X)$ of morphisms $T\to X$, and the elements of
this set are also called \emph{points of $X$ with values in $T$}. If to
every morphism $f:T\to T'$ we associate the map $u'\mapsto u'\circ f$ of $X(T')$
into $X(T)$, we see that, for fixed $X$, $X(T)$ is a \emph{contravariant functor
in $T$}, from the category of preschemes to that of sets. Moreover,
every morphism of preschemes $g:X\to Y$ defines a functorial homomorphism
$X(T)\to Y(T)$, which assigns $g\circ v$ to $v\in X(T)$.
\end{env}

\begin{env}[3.4.2]
\label{I.3.4.2}
Given three sets $P$, $Q$, $R$ and two maps
$\varphi:P\to R$, $\psi:Q\to R$, we give the name \emph{fibre product of $P$
and $Q$ over $R$} (relative to $\varphi$ and $\psi$) to the subset of
\oldpage[I]{112}
the product set $P\times Q$ consisting of the pairs $(p,q)$ such that
$\varphi(p)=\psi(q)$; it is denoted by $P\times_R Q$. Definition
(\hyperref[I.3.2.1]{3.2.1}) of the product of $S$-preschemes can also be
interpreted, with the notation of (\hyperref[I.3.4.1]{3.4.1}), by the
formula
\phantomsection\egalabel{I.3.4.2.1}{3.4.2.1}
\[
  (X\times_S Y)(T)=X(T)\times_{S(T)}Y(T)
  \tag{3.4.2.1}
\]
the maps $X(T)\to S(T)$ and $Y(T)\to S(T)$ corresponding to the structural
morphisms $X\to S$ and $Y\to S$.
\end{env}

\begin{env}[3.4.3]
\label{I.3.4.3}
If we are given a prescheme $S$ and consider only
$S$-preschemes and $S$-morphisms, we shall also denote by $X(T)_S$ the set
$\operatorname{Hom}_S(T,X)$ of $S$-morphisms $T\to X$, allowing ourselves to
omit the subscript $S$ when no confusion is possible; we also say
that the elements of $X(T)_S$ are the \emph{points} (or \emph{$S$-points}
when confusions are to be feared) \emph{of the $S$-prescheme $X$ with values
in the $S$-prescheme $T$}. In particular, an \emph{$S$-section} of $X$
is nothing but a \emph{point of $X$ with values in $S$}. Formula
(\hyperref[I.3.4.2.1]{3.4.2.1}) can then also be written
\phantomsection\egalabel{I.3.4.3.1}{3.4.3.1}
\[
  (X\times_S Y)(T)_S=X(T)_S\times Y(T)_S~;
  \tag{3.4.3.1}
\]
more generally, if $Z$ is an $S$-prescheme and $X$, $Y$, $T$ are
$Z$-preschemes (hence \emph{ipso facto} $S$-preschemes), we have
\phantomsection\egalabel{I.3.4.3.2}{3.4.3.2}
\[
  (X\times_Z Y)(T)_S=X(T)_S\times_{Z(T)_S}Y(T)_S.
  \tag{3.4.3.2}
\]

Note that in order to prove that a triple $(W,r,s)$ consisting of an
$S$-prescheme $W$ and of two $S$-morphisms $r:W\to X$, $s:W\to Y$ is a product
of $X$ and $Y$ (over $Z$), it suffices by definition to verify that for
\emph{every $S$-prescheme $T$}, the diagram
\phantomsection\label{I.3.4.3.product-diagram}
\[
  \xymatrix{
    W(T)_S\ar[r]^{r'}\ar[d]_{s'} & X(T)_S\ar[d]^{\varphi'}\\
    Y(T)_S\ar[r]_{\psi'} & Z(T)_S
  }
\]
makes $W(T)_S$ the fibre product of $X(T)_S$ and $Y(T)_S$ over
$Z(T)_S$, $r'$ and $s'$ corresponding to $r$ and $s$, and $\varphi'$ and
$\psi'$ to the structural morphisms $\varphi:X\to Z$, $\psi:Y\to Z$.
\end{env}

\begin{env}[3.4.4]
\label{I.3.4.4}
When, in the above, $T$ (resp. $S$) is an affine scheme with ring $B$
(resp. $A$), we replace $T$ (resp. $S$) by $B$ (resp. $A$) in the preceding
notation, and we then speak of \emph{points of $X$ with values in the ring
$B$}, or of \emph{points of the $A$-prescheme $X$ with values in the
$A$-algebra $B$}, for the elements of $X(B)$ and of $X(B)_A$ respectively.
Note that $X(B)$ and $X(B)_A$ are now functors \emph{covariant in $B$}. We
shall likewise write $X(T)_A$ for the set of points of the $A$-prescheme $X$
with values in the $A$-prescheme $T$.
\end{env}

\begin{env}[3.4.5]
\label{I.3.4.5}
Consider in particular the case where $T$ is of the form
$\operatorname{Spec}(A)$, where $A$ is a \emph{local} ring; the elements of
$X(A)$ then correspond one-to-one to the \emph{local} homomorphisms
$\mathcal{O}_x\to A$ for $x\in X$ (\hyperref[I.2.4.4]{2.4.4}); we say that
the point $x$ of the underlying space of $X$ is the \emph{locality} of the point
of $X$ with values in $A$ to which it corresponds.

More particularly, we shall call \emph{geometric points} of a prescheme
$X$ the \emph{points of $X$ with values in a field $K$}: giving such a
point therefore amounts to giving its
\oldpage[I]{113}
locality $x$ in the underlying space of $X$ and an \emph{extension} $K$ of
$k(x)$; $K$ will be called the \emph{field of values} of the corresponding
geometric point, and we also say that this geometric point is \emph{localized
at $x$}. We thus define a map $X(K)\to X$, sending a geometric point with
values in $K$ to its locality.

If $S'=\operatorname{Spec}(K)$ is an $S$-prescheme (in other words, if $K$ is
regarded as an extension of a residue field $k(s)$, where $s\in S$) and if $X$
is an $S$-prescheme, an element of $X(K)_S$, or, as one also says, a
\emph{geometric point of $X$ over $s$ with values in $K$}, consists
of giving a $k(s)$-monomorphism of a residue field $k(x)$ into $K$, where
$x$ is a point of $X$ \emph{over $s$} (hence $k(x)$ an extension of
$k(s)$).

More particularly, if $S=\operatorname{Spec}(K)=\{\xi\}$, \emph{the geometric
points of $X$ with values in $K$ are identified with the points $x\in X$ such
that $k(x)=K$}; we also say that the latter are the \emph{points of the
$K$-prescheme $X$ rational over $K$}; if $K'$ is an extension of $K$, the
geometric points of $X$ with values in $K'$ therefore correspond
one-to-one to the points of $X'=X_{(K')}$ rational over $K'$
(\hyperref[I.3.3.14]{3.3.14}).
\end{env}

\begin{lemma}[3.4.6]
\label{I.3.4.6}
Let $X_i$ ($1\leq i\leq n$) be $S$-preschemes, $s$ a point of $S$, and $x_i$
($1\leq i\leq n$) a point of $X_i$ over $s$. Then there exist an
extension $K$ of $k(s)$ and a geometric point of the product
$Y=X_1\times_S X_2\times_S\cdots\times_S X_n$, with values in $K$, whose
projections are localized at the $x_i$.
\end{lemma}

Indeed, there exist $k(s)$-monomorphisms $k(x_i)\to K$ into one and the same
extension $K$ of $k(s)$ (Bourbaki, \emph{Alg.}, Chap. V, \S~4, Prop. 2). The
composites $k(s)\to k(x_i)\to K$ are all identical, so the morphisms
$\operatorname{Spec}(K)\to X_i$ corresponding to $k(x_i)\to K$ are
$S$-morphisms, and we conclude that they define a unique morphism
$\operatorname{Spec}(K)\to Y$. If $y$ is the corresponding point of $Y$, it
is clear that its projection to each of the $X_i$ is $x_i$.

\begin{proposition}[3.4.7]
\label{I.3.4.7}
Let $X_i$ ($1\leq i\leq n$) be $S$-preschemes, and for each index $i$,
let $x_i$ be a point of $X_i$. For there to exist a point $y$ of
$Y=X_1\times_S X_2\times_S\cdots\times_S X_n$ whose projection of index $i$ is
$x_i$ for $1\leq i\leq n$, it is necessary and sufficient that the $x_i$ lie
over one and the same point $s$ of $S$.
\end{proposition}

The condition is obviously necessary; Lemma
(\hyperref[I.3.4.6]{3.4.6}) proves that it is sufficient.

\phantomsection\label{I.3.4.7.underlying-map}
In other words, if $(X)$ denotes the underlying set of $X$, we
see that there is a canonical \emph{surjective} map
$(X\times_S Y)\to (X)\times_{(S)}(Y)$; it should be noted that this map
\emph{is not injective} in general; in other words, \emph{there can exist
several distinct points $z$ in $X\times_S Y$ having the same projections
$x\in X$, $y\in Y$}; this can already be seen when $S$, $X$, $Y$ are
prime spectra of fields $k$, $K$, $K'$, since the tensor product
$K\otimes_k K'$ in general has several distinct prime ideals
(cf. \hyperref[I.3.4.9]{3.4.9}).

\begin{corollary}[3.4.8]
\label{I.3.4.8}
Let $f:X\to Y$ be an $S$-morphism, and $f_{(S')}:X_{(S')}\to Y_{(S')}$ the
$S'$-morphism deduced from $f$ by an extension $S'\to S$ of the base prescheme.
Let $p$ (resp. $q$) be the projection $X_{(S')}\to X$ (resp.
$Y_{(S')}\to Y$); for every subset $M$ of $X$, we have
\phantomsection\label{I.3.4.8.base-change-image}
\[
  q^{-1}(f(M))=f_{(S')}(p^{-1}(M)).
\]
\end{corollary}

\oldpage[I]{114}
Indeed (\hyperref[I.3.3.11]{3.3.11}), $X_{(S')}$ is identified with the
product $X\times_Y Y_{(S')}$ by means of the commutative diagram
\phantomsection\label{I.3.4.8.proof-diagram}
\[
  \xymatrix{
    X\ar[d]_f & X_{(S')}\ar[l]_p\ar[d]^{f_{(S')}}\\
    Y & Y_{(S')}\ar[l]^q
  }
\]

By virtue of (\hyperref[I.3.4.7]{3.4.7}), the relation $q(y')=f(x)$ for
$x\in M$, $y'\in Y_{(S')}$ is equivalent to the existence of an
$x'\in X_{(S')}$ such that $p(x')=x$ and $f_{(S')}(x')=y'$, whence the
corollary.

Lemma (\hyperref[I.3.4.6]{3.4.6}) can be made more precise as follows:

\begin{proposition}[3.4.9]
\label{I.3.4.9}
Let $X$, $Y$ be two $S$-preschemes, $x$ a point of $X$, and $y$ a point of $Y$,
lying over the same point $s\in S$. The set of points of $X\times_S Y$ having
$x$ and $y$ as projections is in canonical one-to-one correspondence with
the set of types of composite extensions of $k(x)$ and $k(y)$, regarded
as extensions of $k(s)$ (Bourbaki, \emph{Alg.}, Chap. VIII, \S~8, Prop. 2).
\end{proposition}

Let $p$ (resp. $q$) be the projection of $X\times_S Y$ to $X$ (resp. $Y$) and
let $E$ be the subspace $p^{-1}(x)\cap q^{-1}(y)$ of the underlying space of
$X\times_S Y$. Note first that the morphisms
$\operatorname{Spec}(k(x))\to S$ and $\operatorname{Spec}(k(y))\to S$
factorize as
$\operatorname{Spec}(k(x))\to\operatorname{Spec}(k(s))\to S$ and
$\operatorname{Spec}(k(y))\to\operatorname{Spec}(k(s))\to S$; since
$\operatorname{Spec}(k(s))\to S$ is a monomorphism
(\hyperref[I.2.4.7]{2.4.7}), it follows from
(\hyperref[I.3.2.4]{3.2.4}) that we have
\phantomsection\label{I.3.4.9.tensor-product-fibre}
\[
  P=\operatorname{Spec}(k(x))\times_S\operatorname{Spec}(k(y))
  =\operatorname{Spec}(k(x))\times_{\operatorname{Spec}(k(s))}
   \operatorname{Spec}(k(y))
  =\operatorname{Spec}\bigl(k(x)\otimes_{k(s)}k(y)\bigr).
\]

We shall define two maps $\alpha:P_0\to E$, $\beta:E\to P_0$,
inverse to each other ($P_0$ denoting the underlying set of the
prescheme $P$). If $i:\operatorname{Spec}(k(x))\to X$ and
$j:\operatorname{Spec}(k(y))\to Y$ are the canonical morphisms
(\hyperref[I.2.4.5]{2.4.5}), we take for $\alpha$ the map of
underlying spaces corresponding to the morphism $i\times_S j$. On the other
hand, every $z\in E$ defines by hypothesis two $k(s)$-monomorphisms
$k(x)\to k(z)$ and $k(y)\to k(z)$, hence a tensor product
$k(s)$-homomorphism $k(x)\otimes_{k(s)}k(y)\to k(z)$ and consequently a
morphism $\operatorname{Spec}(k(z))\to P$; $\beta(z)$ will be the image of $z$
in $P_0$ under this morphism. The verification that $\alpha\circ\beta$ and
$\beta\circ\alpha$ are the identity maps follows from
(\hyperref[I.2.4.5]{2.4.5}) and from the definition of the product
(\hyperref[I.3.2.1]{3.2.1}). Finally, we know that $P_0$ is in one-to-one
correspondence with the set of types of composite extensions of $k(x)$ and
$k(y)$ (Bourbaki, \emph{Alg.}, Chap. VIII, \S~8, Prop. 1).

\subsection{Surjections and injections}
\label{subsection:I.3.5}

\begin{env}[3.5.1]
\label{I.3.5.1}
In general, consider a property $\mathbf{P}$ of morphisms of
preschemes, and the following two propositions:
\begin{enumerate}
  \item[(i)] If $f:X\to X'$, $g:Y\to Y'$ are two $S$-morphisms having
  property $\mathbf{P}$, then $f\times_S g$ has property $\mathbf{P}$.
  \item[(ii)] If $f:X\to Y$ is an $S$-morphism having property
  $\mathbf{P}$, every $S'$-morphism $f_{(S')}:X_{(S')}\to Y_{(S')}$ deduced
  from $f$ by extension of the base prescheme has property
  $\mathbf{P}$.
\end{enumerate}

\oldpage[I]{115}
Since $f_{(S')}=f\times_S 1_{S'}$, we see that if for every prescheme $X$
the \emph{identity} $1_X$ has property $\mathbf{P}$, then (i) implies
(ii); since, on the other hand, $f\times_S g$ is the composite morphism
\phantomsection\label{I.3.5.1.product-composition}
\[
  X\times_S Y\xrightarrow{f\times 1_Y}X'\times_S Y
  \xrightarrow{1_{X'}\times g}X'\times_S Y'
\]
we see that if the \emph{composite} of two morphisms having property
$\mathbf{P}$ also has this property, then (ii) implies (i).
\end{env}

A first application of this remark is the following

\begin{proposition}[3.5.2]
\label{I.3.5.2}
\begin{enumerate}
  \item[(i)] If $f:X\to X'$, $g:Y\to Y'$ are surjective $S$-morphisms,
  then $f\times_S g$ is surjective.
  \item[(ii)] If $f:X\to Y$ is a surjective $S$-morphism, then $f_{(S')}$ is
  surjective for every extension $S'$ of the base prescheme.
\end{enumerate}
\end{proposition}

Since the composite of two surjections is a surjection, it suffices to prove
(ii); but this proposition follows at once from
(\hyperref[I.3.4.8]{3.4.8}) applied to $M=X$.

\begin{proposition}[3.5.3]
\label{I.3.5.3}
For a morphism $f:X\to Y$ to be surjective, it is necessary and sufficient
that for every field $K$ and every morphism $\operatorname{Spec}(K)\to Y$,
there exist an extension $K'$ of $K$ and a morphism
$\operatorname{Spec}(K')\to X$ making the diagram
\phantomsection\label{I.3.5.3.surjectivity-diagram}
\[
  \xymatrix{
    X\ar[d]_f & \operatorname{Spec}(K')\ar[l]\ar[d]\\
    Y & \operatorname{Spec}(K)\ar[l]
  }
\]
commutative.
\end{proposition}

The condition is sufficient, since for every $y\in Y$ it suffices to apply it
to a morphism $\operatorname{Spec}(K)\to Y$ corresponding to a monomorphism
$k(y)\to K$, $K$ being an extension of $k(y)$
(\hyperref[I.2.4.6]{2.4.6}). Conversely, suppose $f$ surjective, and let
$y\in Y$ be the image of the unique point of $\operatorname{Spec}(K)$; there
exists $x\in X$ such that $f(x)=y$; consider the corresponding monomorphism
$k(y)\to k(x)$ (\hyperref[I.2.2.1]{2.2.1}); it then suffices to take for
$K'$ an extension of $k(y)$ such that there exist $k(y)$-monomorphisms of
$k(x)$ and of $K$ into $K'$ (Bourbaki, \emph{Alg.}, Chap. V, \S~4, Prop. 2);
the morphism $\operatorname{Spec}(K')\to X$ corresponding to $k(x)\to K'$
answers the question.

In the language introduced in (\hyperref[I.3.4.5]{3.4.5}), we can also say
that \emph{every geometric point of $Y$ with values in $K$ comes
from a geometric point of $X$ with values in an extension of $K$}.

\begin{definition}[3.5.4]
\label{I.3.5.4}
A morphism of preschemes $f:X\to Y$ is said to be \emph{universally
injective}, or a \emph{radicial morphism}, if for every field $K$,
the corresponding map $X(K)\to Y(K)$ is injective.
\end{definition}

It follows at once from the definitions that every \emph{monomorphism of
preschemes} (T, 1.1) is radicial.

\begin{env}[3.5.5]
\label{I.3.5.5}
For a morphism $f:X\to Y$ to be radicial, it suffices that the condition of
Definition (\hyperref[I.3.5.4]{3.5.4}) be satisfied for every
\emph{algebraically closed} field. Indeed, if $K$ is an arbitrary field and
$K'$ an algebraically closed extension of $K$, the diagram
\phantomsection\label{I.3.5.5.algebraic-closure-diagram}
\[
  \xymatrix{
    X(K)\ar[r]^\alpha\ar[d]_\varphi & Y(K)\ar[d]^{\varphi'}\\
    X(K')\ar[r]_{\alpha'} & Y(K')
  }
\]

\oldpage[I]{116}
is commutative, $\varphi$ and $\varphi'$ coming from the morphism
$\operatorname{Spec}(K')\to\operatorname{Spec}(K)$, and $\alpha$ and $\alpha'$
corresponding to $f$. Now, $\varphi$ is injective, and so is
$\alpha'$ by hypothesis; hence $\alpha$ is necessarily injective.
\end{env}

\begin{proposition}[3.5.6]
\label{I.3.5.6}
Let $f:X\to Y$, $g:Y\to Z$ be two morphisms of preschemes.
\begin{enumerate}
  \item[(i)] If $f$ and $g$ are radicial, so is $g\circ f$.
  \item[(ii)] Conversely, if $g\circ f$ is radicial, so is
  $f$.
\end{enumerate}
\end{proposition}

Taking into account Definition (\hyperref[I.3.5.4]{3.5.4}), the proposition
reduces to the corresponding assertions for the maps
$X(K)\to Y(K)\to Z(K)$, which are obvious.

\begin{proposition}[3.5.7]
\label{I.3.5.7}
\begin{enumerate}
  \item[(i)] If the $S$-morphisms $f:X\to X'$, $g:Y\to Y'$ are radicial,
  so is $f\times_S g$.
  \item[(ii)] If the $S$-morphism $f:X\to Y$ is radicial, so is
  $f_{(S')}:X_{(S')}\to Y_{(S')}$ for every extension $S'\to S$ of the
  base prescheme.
\end{enumerate}
\end{proposition}

In view of (\hyperref[I.3.5.1]{3.5.1}), it suffices to prove (i). We have
seen (\hyperref[I.3.4.2.1]{3.4.2.1}) that
\phantomsection\label{I.3.5.7.field-valued-product}
\[
  (X\times_S Y)(K)=X(K)\times_{S(K)}Y(K),\qquad
  (X'\times_S Y')(K)=X'(K)\times_{S(K)}Y'(K)
\]
the map $(X\times_S Y)(K)\to(X'\times_S Y')(K)$ corresponding to
$f\times_S g$ is then identified with $(u,v)\mapsto(f\circ u,g\circ v)$, and
the proposition follows at once.

\begin{proposition}[3.5.8]
\label{I.3.5.8}
For a morphism $f=(\psi,\theta):X\to Y$ to be radicial, it is necessary and
sufficient that $\psi$ be injective and that, for every $x\in X$, the
monomorphism $\theta^x:k(\psi(x))\to k(x)$ make $k(x)$ a radicial extension of
$k(\psi(x))$.
\end{proposition}

Suppose $f$ radicial, and let us first show that the relation
$\psi(x_1)=\psi(x_2)=y$ necessarily implies $x_1=x_2$. Indeed, there
exist a field $K$, an extension of $k(y)$, and $k(y)$-monomorphisms
$k(x_1)\to K$, $k(x_2)\to K$ (Bourbaki, \emph{Alg.}, Chap. V, \S~4,
Prop. 2); the corresponding morphisms
$u_1:\operatorname{Spec}(K)\to X$, $u_2:\operatorname{Spec}(K)\to X$ are
therefore such that $f\circ u_1=f\circ u_2$, hence $u_1=u_2$ by hypothesis,
and this implies in particular $x_1=x_2$. Now regard $k(x)$ as an
extension of $k(\psi(x))$ by means of $\theta^x$: if $k(x)$ is not a
radicial extension of $k(\psi(x))$, there exist two distinct
$k(\psi(x))$-monomorphisms of $k(x)$ into an algebraically
closed extension $K$ of $k(\psi(x))$, and the two corresponding morphisms
$\operatorname{Spec}(K)\to X$ would violate the hypothesis. Conversely, taking
into account (\hyperref[I.2.4.6]{2.4.6}), it is immediate that the
conditions of the statement are sufficient for $f$ to be radicial.

\begin{corollary}[3.5.9]
\label{I.3.5.9}
If $A$ is a ring and $S$ a multiplicative subset of $A$, the canonical
morphism $\operatorname{Spec}(S^{-1}A)\to\operatorname{Spec}(A)$ is
radicial.
\end{corollary}

Indeed, this morphism is a monomorphism
(\hyperref[I.1.6.2]{1.6.2}).

\begin{corollary}[3.5.10]
\label{I.3.5.10}
Let $f:X\to Y$ be a radicial morphism, $g:Y'\to Y$ a morphism, and let
$X'=X_{(Y')}=X\times_Y Y'$. Then the radicial morphism $f_{(Y')}$
(\hyperref[I.3.5.7]{3.5.7 (ii)}) is a bijection of the underlying space
$X'$ onto $g^{-1}(f(X))$; moreover, for every field $K$, the set $X'(K)$
is identified with the subset of $Y'(K)$ which is the inverse image, under the
map $Y'(K)\to Y(K)$ (corresponding to $g$), of the subset $X(K)$ of $Y(K)$.
\end{corollary}

The first assertion follows at once from
(\hyperref[I.3.5.8]{3.5.8}) and from
(\hyperref[I.3.4.8]{3.4.8}); the second, from the commutativity of the
diagram
\oldpage[I]{117}
\phantomsection\label{I.3.5.10.base-change-points-diagram}
\[
  \xymatrix{
    X'(K)\ar[r]\ar[d] & Y'(K)\ar[d]\\
    X(K)\ar[r] & Y(K)
  }
\]

\begin{remark}[3.5.11]
\label{I.3.5.11}
We shall say that a morphism $f=(\psi,\theta)$ of preschemes is
\emph{injective} if the map $\psi$ is injective. For a morphism
$f=(\psi,\theta):X\to Y$ to be radicial, it is necessary and sufficient that
for every morphism $Y'\to Y$, the morphism $f_{(Y')}:X_{(Y')}\to Y'$ be
injective (which justifies the terminology of \emph{universally injective}
morphism). Indeed, the condition is necessary by virtue of
(\hyperref[I.3.5.7]{3.5.7, (ii)}) and of
(\hyperref[I.3.5.8]{3.5.8}). Conversely, the condition first implies that
$\psi$ is injective; if for some $x\in X$ the monomorphism
$\theta^x:k(\psi(x))\to k(x)$ were not radicial, there would be an extension
$K$ of $k(\psi(x))$ and two distinct morphisms
$\operatorname{Spec}(K)\to X$ corresponding to the same morphism
$\operatorname{Spec}(K)\to Y$ (\hyperref[I.3.5.8]{3.5.8}). But then, on
setting $Y'=\operatorname{Spec}(K)$, there would be two distinct $Y'$-sections
of $X_{(Y')}$ (\hyperref[I.3.3.14]{3.3.14}), which contradicts the
hypothesis that $f_{(Y')}$ is injective.
\end{remark}

\subsection{Fibres.}
\label{subsection:I.3.6}

\begin{proposition}[3.6.1]
\label{I.3.6.1}
Let $f:X\to Y$ be a morphism, $y$ a point of $Y$, and $\mathfrak{a}_y$ an
ideal of definition of $\mathcal{O}_y$ for the
$\mathfrak{m}_y$-preadic topology. The projection
$p:X\times_Y\operatorname{Spec}(\mathcal{O}_y/\mathfrak{a}_y)\to X$ is a
homeomorphism of the underlying space of the prescheme
$X\times_Y\operatorname{Spec}(\mathcal{O}_y/\mathfrak{a}_y)$ onto the fibre
$f^{-1}(y)$ endowed with the topology induced by that of the underlying space
of $X$.
\end{proposition}

Since $\operatorname{Spec}(\mathcal{O}_y/\mathfrak{a}_y)\to Y$ is radicial
(\hyperref[I.3.5.4]{3.5.4} and \hyperref[I.2.4.7]{2.4.7}) and
$\operatorname{Spec}(\mathcal{O}_y/\mathfrak{a}_y)$ is reduced to a single
point, the ideal $\mathfrak{m}_y/\mathfrak{a}_y$ being nilpotent by hypothesis
(\hyperref[I.1.1.12]{1.1.12}), we already know
(\hyperref[I.3.5.10]{3.5.10} and \hyperref[I.3.3.4]{3.3.4}) that $p$
identifies, \emph{as a set}, the underlying space of
$X\times_Y\operatorname{Spec}(\mathcal{O}_y/\mathfrak{a}_y)$ with $f^{-1}(y)$;
everything comes down to proving that $p$ is a homeomorphism. By virtue of
(\hyperref[I.3.2.7]{3.2.7}), the question is local on $X$ and $Y$, and we
may therefore assume that $X=\operatorname{Spec}(B)$,
$Y=\operatorname{Spec}(A)$, $B$ being an $A$-algebra. The morphism $p$
then corresponds to the homomorphism
$1\otimes\varphi:B\to B\otimes_A A'$, where $A'=A_y/\mathfrak{a}_y$ and
$\varphi$ is the canonical map of $A$ into $A'$. Now, every element of
$B\otimes_A A'$ can be written
\phantomsection\label{I.3.6.1.localization-fraction}
\[
  \sum_i b_i\otimes\varphi(a_i)/\varphi(s)
  =\left(\sum_i(a_i b_i\otimes 1)\right)(1\otimes\varphi(s))^{-1},
\]
where $s\notin\mathfrak{j}_y$, and Proposition
(\hyperref[I.1.2.4]{1.2.4}) is applicable.

\begin{env}[3.6.2]
\label{I.3.6.2}
Throughout the rest of this Treatise, when we regard a fibre
$f^{-1}(y)$ of a morphism as endowed with a structure of $k(y)$-prescheme,
\emph{it will always be the prescheme obtained by transporting the structure
of $X\times_Y\operatorname{Spec}(k(y))$ by the projection to $X$}. We
shall also write this last product as $X\otimes_Y k(y)$, or
$X\otimes_{\mathcal{O}_Y}k(y)$; more generally, if $B$ is an
$\mathcal{O}_y$-algebra, we shall denote by $X\otimes_Y B$ or
$X\otimes_{\mathcal{O}_Y}B$ the product
$X\times_Y\operatorname{Spec}(B)$.
\end{env}

With the preceding convention, it follows from
(\hyperref[I.3.5.10]{3.5.10}) that the points of $X$ with values in an
extension $K$ of $k(y)$ are identified with the \emph{points of $f^{-1}(y)$
with values in $K$}.

\begin{env}[3.6.3]
\label{I.3.6.3}
Let $f:X\to Y$, $g:Y\to Z$ be two morphisms, and $h=g\circ f$ their composite;
for every $z\in Z$, the fibre $h^{-1}(z)$ is a prescheme isomorphic to
\phantomsection\label{I.3.6.3.fibre-composition}
\[
  X\times_Z\operatorname{Spec}(k(z))
  =(X\times_Y Y)\times_Z\operatorname{Spec}(k(z))
  =X\times_Y g^{-1}(z)
\]
\oldpage[I]{118}
In particular, if $U$ is an open subset of $X$, the prescheme induced on
$U\cap f^{-1}(y)$ by the prescheme $f^{-1}(y)$ is isomorphic to
$f_U^{-1}(y)$ ($f_U$ being the restriction of $f$ to $U$).
\end{env}

\begin{proposition}[3.6.4][``transitivity of fibres'']
\label{I.3.6.4}
Let $f:X\to Y$, $g:Y'\to Y$ be two morphisms; set
$X'=X\times_Y Y'=X_{(Y')}$, and
$f'=f_{(Y')}:X'\to Y'$. For every $y'\in Y'$, if we set $y=g(y')$, the
prescheme $f'^{-1}(y')$ is isomorphic to
$f^{-1}(y)\otimes_{k(y)}k(y')$.
\end{proposition}

Indeed, this amounts to noting that the two preschemes
$(X\otimes_Y k(y))\otimes_{k(y)}k(y')$ and
$(X\times_Y Y')\otimes_{Y'}k(y')$ are both canonically isomorphic to
$X\times_Y\operatorname{Spec}(k(y'))$ by
(\hyperref[I.3.3.9.1]{3.3.9.1}).

In particular, if $V$ is an open neighbourhood of $y$ in $Y$, and if
$f_V$ denotes the restriction of $f$ to the prescheme induced on $f^{-1}(V)$,
the preschemes $f^{-1}(y)$ and $f_V^{-1}(y)$ are canonically identified.

\begin{proposition}[3.6.5]
\label{I.3.6.5}
Let $f:X\to Y$ be a morphism, $y$ a point of $Y$, $Z$ the local prescheme
$\operatorname{Spec}(\mathcal{O}_y)$, and $p=(\psi,\theta)$ the projection
$X\times_Y Z\to X$; then $p$ is a homeomorphism of the underlying space of
$X\times_Y Z$ onto the subspace $f^{-1}(Z)$ of $X$ (when the underlying space
of $Z$ is identified with a subspace of $Y$, cf.
(\hyperref[I.2.4.2]{2.4.2})), and for every $t\in X\times_Y Z$, if we set
$z=\psi(t)$, $\theta_t^\sharp$ is an isomorphism of $\mathcal{O}_z$ onto
$\mathcal{O}_t$.
\end{proposition}

Since $Z$ (identified with a subspace of $Y$) is contained in every affine
open set containing $y$ (\hyperref[I.2.4.2]{2.4.2}), we can, as in
(\hyperref[I.3.6.1]{3.6.1}), reduce to the case where
$X=\operatorname{Spec}(A)$ and $Y=\operatorname{Spec}(B)$ are affine schemes,
$A$ being a $B$-algebra. Then $X\times_Y Z$ is the prime spectrum
of $A\otimes_B B_y$, and this ring is canonically identified with $S^{-1}A$,
where $S$ is the image of $B-\mathfrak{j}_y$ in $A$ (0, 1.5.2); since $p$
then corresponds to the canonical homomorphism $A\to S^{-1}A$, the proposition
follows from (\hyperref[I.1.6.2]{1.6.2}).

\subsection{Application: reduction of a prescheme mod.
\texorpdfstring{$\mathfrak{J}$\protect\footnotemark}{J}}
\label{subsection:I.3.7}
\label{I.3.7}
\footnotetext{This subsection, which makes use of notions and results from
the rest of Chapter~I and from Chapter~II, will not be used in the sequel,
and is intended only for readers familiar with classical algebraic geometry.}

\begin{env}[3.7.1]
\label{I.3.7.1}
Let $A$ be a ring, $X$ an $A$-prescheme, and $\mathfrak{J}$ an ideal of $A$;
then $X_0=X\otimes_A(A/\mathfrak{J})$ is an
$(A/\mathfrak{J})$-prescheme, which is sometimes said to be deduced from $X$ by
\emph{reduction} mod. $\mathfrak{J}$.
\end{env}

\begin{env}[3.7.2]
\label{I.3.7.2}
This terminology is used above all when $A$ is a \emph{local
ring} and $\mathfrak{J}$ its maximal ideal, so that $X_0$ is a
prescheme over the residue field $k=A/\mathfrak{J}$ of $A$.

When moreover $A$ is an integral domain, with field of fractions $K$, one can
also consider the $K$-prescheme $X'=X\otimes_A K$. By an abuse of language which
we shall not use, it was customary until now to say that $X_0$ is
\emph{deduced from $X'$} by reduction mod. $\mathfrak{J}$. In the cases where
this language was used, $A$ was a local ring of dimension 1 (most
often a discrete valuation ring), and it was understood (more or less
explicitly) that the given $K$-prescheme $X'$ was a
closed sub-prescheme of a $K$-prescheme $P'$ (in fact, a standard projective
space $\mathbf{P}_K^r$, cf. II, 4.1.1), itself of the form
$P'=P\otimes_A K$, where $P$ is a given $A$-prescheme (in fact, the $A$-scheme
$\mathbf{P}_A^r$, with the notation of II, 4.1.1). In our language, the
definition of $X_0$ in terms of $X'$ is formulated as follows:

Consider the affine scheme $Y=\operatorname{Spec}(A)$, consisting of two
points, the unique
\oldpage[I]{119}
closed point $y=\mathfrak{J}$ and the generic point $(0)$, the set $U$ reduced
to the generic point therefore being an open set $U=\operatorname{Spec}(K)$ in
$Y$. If $X$ is an $A$-prescheme (in other words a $Y$-prescheme),
$X\otimes_A K=X'$ is nothing but the prescheme induced by $X$ on
$\psi^{-1}(U)$, where $\psi$ denotes the structural morphism $X\to Y$. In
particular, if $\varphi$ is the structural morphism $P\to Y$, a
closed sub-prescheme $X'$ of $P'=\varphi^{-1}(U)$ is therefore a (locally
closed) sub-prescheme of $P$. If $P$ is Noetherian (for example if $A$ is
Noetherian and $P$ is of finite type over $A$), there exists a smallest
closed sub-prescheme $X=\overline{X'}$ of $P$ which majorizes $X'$ (9.5.10),
and $X'$ is the prescheme induced by $X$ on the open set
$\varphi^{-1}(U)\cap X$, hence is isomorphic to $X\otimes_A K$ (9.5.10).
\emph{The immersion of $X'$ into $P'=P\otimes_A K$ therefore makes it possible
to regard $X'$ canonically as being of the form $X'=X\otimes_A K$, where
$X$ is an $A$-prescheme.} One can then consider the prescheme reduced mod.
$\mathfrak{J}$, $X_0=X\otimes_A k$, which is moreover nothing but the fibre
$\psi^{-1}(y)$ of the closed point $y$. Until now, for lack of an adequate
terminology, one had avoided introducing the $A$-prescheme $X$ explicitly. It
should however be noted that all the assertions usually made
about the prescheme ``reduced mod. $\mathfrak{J}$'' $X_0$ must be regarded
as consequences of more complete assertions concerning $X$ itself, and
can be formulated and understood satisfactorily only by
interpreting them in this way. It seems moreover that the hypotheses made always
amount to hypotheses on $X$ itself (independently of an immersion of $X'$ into
a $\mathbf{P}_K^r$ being given beforehand), which makes it possible to
give more intrinsic statements.
\end{env}

\begin{env}[3.7.3]
\label{I.3.7.3}
Finally, let us point out a very special fact, which has no doubt contributed
to delaying the conceptual clarification of the situation considered here: if
$A$ is a discrete valuation ring and $X$ is \emph{proper over $A$}
(which is in fact the case if $X$ is a closed sub-prescheme of a
$\mathbf{P}_A^r$, cf. II, 5.5.4), the points of $X$ with values in $A$ and the
points of $X'$ with values in $K$ are in one-to-one correspondence
(II, 7.3.8). This is why one has often believed oneself to be proving results
concerning $X'$, whereas in reality one was proving statements concerning
$X$, which remain valid (in this form) when the local base ring is no longer
assumed to be of dimension 1.
\end{env}
